\section{LIMITATIONS}

The frozen domain test set contains 170 passages and lacks full mutually blind
double annotation with formal adjudication. A fresh two-human non-test
calibration reached exact micro-F1 of 0.8000, meeting the original minimum, but
not the later conservative 0.85 target. It cannot substitute for frozen-test
formal IAA, and \entitytype{TECH\_SCHEME} F1 of 0.153846 remains the principal
dataset-quality limitation.

The evaluated NER scope comprises A/B/C/D at $K=5$ over three seeds: 12 domain
run units and 12 CLUENER run units over all 1,343 labeled official development
rows. Additional public benchmarks, $K=1$/$K=10$ scans, additional seeds, a
supervised competitive baseline, and per-type uncertainty analysis remain
outside this scope. Means and sample standard deviations describe observed run
variation rather than statistical significance or broad cross-dataset and
$K$-robustness evidence.

The two three-seed blocks are NER-only. A separate four-cell seed-42 Stage-R
audit reports exact RE F1 of 0.040541/0.056972/0.057971/0.045381 for A/B/C/D,
with zero \texttt{HAS\_METRIC} true positives and D below C. These low scores
bound positive RE, RAG-for-RE, and CoVe claims while preserving an auditable
two-stage execution record. The official CLUENER test file has no usable
offline gold; the public result therefore uses labeled official development
rows and is not an official-test score.

The CLUENER block records 138/16,116 (0.856292\%) mechanical parse failures
after the fixed retry ladder. Its 30-record support bank, 61-entry
support-derived lexicon, and labels-stripped RAG replay are schema-specific
resources rather than transferred equipment-domain assets. Consequently, the
block supports limited descriptive directional replication, not broad
portability.

The original frozen inventory contains 60 planned final model identities. Four
domain Stage-N-plus-Stage-R cells and 12 CLUENER NER cells are complete; eight
domain identities retain Stage N only and 36 are absent, giving literal
accounting of 16/60. The other 44 identities are route-closed by the human
decision, not failed experiments. This boundary precludes a complete-matrix or
end-to-end IE claim.

Finally, false-positive entity count is an entity-level extraction-error
quantity rather than an open-domain factuality metric. The qualitative C/D
audit uses seed 42, and D introduces errors as well as removing them. EC-FSIE
was evaluated with one local extractor family and one locked
embedding/reranking stack; the dependency lock and trace artifacts support
reproduction of this implementation, while operating-point behavior across
other model families and deployment stacks remains to be established.
